Deployed AI systems change through instructions, skills, memories, tools, and orchestration long before their model weights are updated. These edits accumulate experience, but their persistence alone does not establish that the system is improving. This survey examines how 256 production and research systems turn experience into lasting edits outside the weights. We distinguish two decisions: whether an edit improves the deployment that produced it, and whether enough evidence supports sharing it more widely. These decisions define a local improvement loop and a promotion loop. The feedback these loops receive (exact checks, observed outcomes, judgments, or delayed outcomes) supports different acceptance rules and exposes different limitations. In every case, an accepted edit is supported only for the measure and tasks on which it was evaluated. We use this observation to develop five conditions for artifact-layer optimization, a classification of feedback and gate types, and four properties for auditing promotion across organizations. The resulting framework identifies what existing systems can already automate and where new evidence or human judgment is still required. Nearly every recently added system with a local loop has an automatic acceptance test (89 of 96), while failure records and credit assignment remain less developed. Promotion is more fragmented: the surveyed systems provide evaluation, shared artifact pools, lineage, access control, and rollback in different combinations, but none documents their integration into governed promotion across organizations. The feedback available to a deployment helps explain this pattern.
